% Part: many-valued-logic
% Chapter: three-valued-logics
% Section: kleene

\documentclass[../../../include/open-logic-section]{subfiles}

\begin{document}

\olfileid{mvl}{thr}{skl}

\olsection{کلینی دیاں منطقاں}

سٹیفن کلینی نے تِن قدراں والیاں دو منطقاں پیش کیتیاں۔ اوہناں
دا خیال اوس منطق توں آیا جیہدے وچ صداقتی قدراں نوں حسابی
طریقیاں دے نتیجے سمجھدے آں۔ کوئی طریقہ $\True$ یا $\False$
دے سکدا اے، پر ایہ وی ہو سکدا اے کہ اوہ مُکے ای نہ۔ اوہ
صورت وچ متعلقہ صداقتی قدر غیر معین ہوندی اے، جیہنوں~$\Undef$
نال لکھدے آں۔

قضیے~$!A$ دی نفی کڈھن لئی پہلاں~$!A$ دی قدر کڈھو، تے
فیر نتیجے دی اُلٹ قدر واپس دیو۔ جے $!A$ دا حساب نہیں مُکدا،
تاں سارا طریقہ وی نہیں مُکدا۔ سو $\Undef$ دی نفی $\Undef$ اے۔

عطف $!A \land !B$ دا حساب کرن دے دو طریقے نیں۔ اک ایہ اے
کہ پہلاں~$!A$ دا حساب کرو، فیر $!B$ دا۔ عطف دا نتیجہ $\True$
ہووے گا جے دوہاں دا نتیجہ~$\True$ ہووے۔ ہور صورت وچ $\False$
ہووے گا۔ جے دوہاں وچوں کوئی حساب نہ رُکے، تاں سارا طریقہ وی
نہیں رُکدا۔ سو ایس صورت وچ جے اک جزو غیر معین ہووے، تاں عطف
وی غیر معین اے۔ فصل بارے وی اوہی گل اے۔

پر جے اسیں $!A$ تے $!B$ دیاں قدراں متوازی کڈھ سکیے، تاں
ہور سوکھ اے۔ جے دو طریقیاں وچوں اک رُک کے $\False$ واپس
دے دے، تاں اسیں رک سکدے آں، کیوں جے جواب لازماً جھوٹ اے۔
سو ایس صورت وچ اک جزو جھوٹ ہووے تاں عطف جھوٹ اے، بھاویں
دوجا جزو غیر معین ہووے۔ ایسے طرح فصل دا حساب متوازی کریے،
تاں جدوں دوہاں جزواں وچوں اک دا طریقہ سچ واپس دے دے، اسیں
رک سکدے آں۔ فیر فصل لازماً سچ اے۔ سو ایس صورت وچ مرکب
قضیے دا نتیجہ پتا لگ سکدا اے، بھاویں اوہدے جزواں وچوں اک
غیر معین ہووے۔ ایس سمجھ مطابق $\Undef$ نوں «غیر معین» دی
تھاں «نامعلوم» پڑھ سکدے آں۔

ایہ دو سمجھاں توں کلینی دی مضبوط تے کمزور منطق بندی اے۔
شرطیہ نوں $\lnot !A \lor !B$ دے ہم ارز متعین کردے آں۔

\begin{defn}
\emph{کلینی دی مضبوط منطق}~$\LogKs$ ایس میٹرکس نال متعین اے۔
\begin{enumerate}
  \item معیاری قضیاتی زبان $\Lang L_0$ جیہدے وچ
  $\lnot$، $\land$، $\lor$، $\lif$ نیں۔
  \item صداقتی قدراں دا سیٹ $V = \{\True, \Undef, \False\}$۔
  \item $\True$ اکّو مقرر قدر اے، یعنی $V^+ = \{\True\}$۔
  \item صداقتی فنکشن ہِٹھلے جدولاں توں ملدے نیں۔
  \begin{center}
    \begin{tabular}{c|c} 
      $\tf{\lnot}$ & \\ 
      \hline  
      $\True$ & $\False$ \\ 
      $\Undef$ & $\Undef$ \\
      $\False$ & $\True$ 
    \end{tabular}
    \quad
    \begin{tabular}{c|ccc} 
      $\tf{\land}[\LogKs]$ & $\True$ & $\Undef$ & $\False$ \\ 
      \hline 
      $\True$ & $\True$ & $\Undef$ & $\False$ \\ 
      $\Undef$ & $\Undef$ & $\Undef$ & $\False$\\ 
      $\False$ & $\False$ & $\False$ & $\False$ 
    \end{tabular}
    \\[2ex]
    \begin{tabular}{c|ccc} 
      $\tf{\lor}[\LogKs]$ & $\True$ & $\Undef$ & $\False$ \\ 
      \hline 
      $\True$ & $\True$ & $\True$ & $\True$ \\ 
      $\Undef$ & $\True$ & $\Undef$ & $\Undef$ \\
      $\False$ & $\True$ & $\Undef$ & $\False$ 
    \end{tabular}
    \quad
    \begin{tabular}{c|ccc} 
      $\tf{\lif}[\LogKs]$ & $\True$ & $\Undef$ & $\False$ \\ 
      \hline 
      $\True$ & $\True$ & $\Undef$ & $\False$ \\ 
      $\Undef$ & $\True$ & $\Undef$ & $\Undef$  \\ 
      $\False$ & $\True$ & $\True$ & $\True$ 
    \end{tabular}
  \end{center}
\end{enumerate}
\end{defn}

\begin{defn}
\emph{کلینی دی کمزور منطق}~$\LogKw$ ایس میٹرکس نال متعین اے۔
\begin{enumerate}
  \item معیاری قضیاتی زبان $\Lang L_0$ جیہدے وچ
  $\lnot$، $\land$، $\lor$، $\lif$ نیں۔
  \item صداقتی قدراں دا سیٹ $V = \{\True, \Undef, \False\}$۔
  \item $\True$ اکّو مقرر قدر اے، یعنی $V^+ = \{\True\}$۔
  \item صداقتی فنکشن ہِٹھلے جدولاں توں ملدے نیں۔
  \begin{center}
    \begin{tabular}{c|c} 
      $\tf{\lnot}$ & \\ 
      \hline  
      $\True$ & $\False$ \\ 
      $\Undef$ & $\Undef$ \\
      $\False$ & $\True$ 
    \end{tabular}
    \quad
    \begin{tabular}{c|ccc} 
      $\tf{\land}[\LogKw]$ & $\True$ & $\Undef$ & $\False$ \\ 
      \hline 
      $\True$ & $\True$ & $\Undef$ & $\False$ \\ 
      $\Undef$ & $\Undef$ & $\Undef$ & $\Undef$\\ 
      $\False$ & $\False$ & $\Undef$ & $\False$ 
    \end{tabular}
    \\[2ex]
    \begin{tabular}{c|ccc} 
      $\tf{\lor}[\LogKw]$ & $\True$ & $\Undef$ & $\False$ \\ 
      \hline 
      $\True$ & $\True$ & $\Undef$ & $\True$ \\ 
      $\Undef$ & $\Undef$ & $\Undef$ & $\Undef$ \\
      $\False$ & $\True$ & $\Undef$ & $\False$ 
    \end{tabular}
    \quad
    \begin{tabular}{c|ccc} 
      $\tf{\lif}[\LogKw]$ & $\True$ & $\Undef$ & $\False$ \\ 
      \hline 
      $\True$ & $\True$ & $\Undef$ & $\False$ \\ 
      $\Undef$ & $\Undef$ & $\Undef$ & $\Undef$  \\ 
      $\False$ & $\True$ & $\Undef$ & $\True$ 
    \end{tabular}
  \end{center}
\end{enumerate}
\end{defn}

\begin{prop}
  $\LogKs$ تے $\LogKw$ وچ کوئی ہمیشہ صادق فارمولا نہیں۔
\end{prop}

\begin{proof}
  جے $\pAssign v(p) = \Undef$ سارے !!{propositional variable}s~$p$
  لئی، تاں ہر فارمولا~$!A$ دی صداقتی قدر~$\pValue v(!A) =
  \Undef$ ہووے گی، کیوں جے دوہاں منطقاں وچ
  \[
    \tf{\lnot}(\Undef) = \tf{\lor}(\Undef, \Undef) = \tf{\land}(\Undef,
  \Undef) = \tf{\lif}(\Undef, \Undef) = \Undef
  \]
  اے۔ کیوں جے $\Undef \notin V^+$، چاہے منطق $\LogKs$ ہووے یا
  $\LogKw$، ایس لئی ایس !!{valuation} اُتے $!A$ دی قدر مقرر نہیں۔
\end{proof}

بھاویں کلینی دیاں کمزور تے مضبوط دوہاں منطقاں وچ کوئی
ہمیشہ صادق فارمولا نہیں، اوہناں دے منطقی نتیجے دے تعلق
معمولی نہیں۔

\begin{prob}
  ہِٹھلے تعلقاں وچوں کیہڑے (الف) کلینی دی مضبوط منطق وچ تے
  (ب) اوہدی کمزور منطق وچ قائم نیں؟ ہر اک لئی صداقتی جدول دیو۔
  \begin{enumerate}
    \item $p, p \lif q \Entails q$
    \item $p \lor q, \lnot p \Entails q$
    \item $p \land q \Entails p$
    \item $p \Entails p \land p$
    \item $p \Entails p \lor q$
  \end{enumerate}
\end{prob}

دمتری بوچوار نے $\Undef$ نوں «بے معنی» سمجھیا۔ اوہنے جھوٹے
دے تناقض ورگے تناقض حل کرن دی کوشش وچ مقرر کیتا کہ تناقض
والے جملیاں دی قدر~$\Undef$ ہووے۔ اوہنے اک منطق پیش کیتی جیہڑی
بنیادی طور اُتے کلینی دی کمزور منطق اے، جیہدے وچ کجھ ہور
رابطے شامل نیں۔ اوہناں وچوں دو «باہرلی نفی» تے «غیر معین اے»
والا عامل نیں۔
\begin{center}
  \begin{tabular}{c|c} 
    $\tf{\sim}$ & \\ 
    \hline  
    $\True$ & $\False$ \\ 
    $\Undef$ & $\True$ \\
    $\False$ & $\True$ 
  \end{tabular}
\quad
  \begin{tabular}{c|c} 
    $\tf{+}$ & \\ 
    \hline  
    $\True$ & $\False$ \\ 
    $\Undef$ & $\True$ \\
    $\False$ & $\False$ 
  \end{tabular}
\end{center}

\begin{prob}
  کی تُسیں بوچوار دی منطق وچ $\sim$ نوں $\lnot$ تے~$+$ نال
  متعین کر سکدے او؟ یعنی !!a{formula} لبھو جیہدے وچ بس
  !!{propositional
  variable}~$p$ ہووے تے $\sim$ نہ آوے، تے جیہدی صداقتی قدر
  ہمیشہ~$\mathord{\sim}p$ دی قدر دے برابر ہووے۔ اپنی گل صحیح
  ثابت کرن لئی صداقتی جدول دیو۔
\end{prob}

\end{document}
